逆函数定理常以下述形式使用, 称为隐函数定理.

\begin{Theorem}{隐函数定理}
\label{thm:banach-implicit-function}
设 $X$, $Y$ 和 $Z$ 为 Banach 空间, $U$ 和 $V$ 分别为 $X$ 和 $Y$ 中原点的邻域. 设 $F$ 从 $U\times V$ 到 $Z$ 连续可微, 并假设 $D_yF(0,0)\in\mathcal{L}(Y,Z)$ 既单又满. 再假设 $F(0,0)=0$. 则存在 $X$ 中原点的邻域 $W$ 以及连续可微映射 $f:W\to Y$, 使 $F(x,f(x))=0$. 此外, 对足够小的 $x$ 和 $y$, $f(x)$ 是方程 $F(x,y)=0$ 的唯一解 $y$.
\end{Theorem}

\begingroup
\renewcommand{\endProof}{\par\nobreak\hfill$\blacksquare$\par\vspace{0.1cm}}
\begin{Proof}
为证明该定理, 只需考虑由 $\widetilde F(x,y)=(x,F(x,y))$ 定义的映射 $\widetilde F:U\times V\to X\times Z$, 并应用前一定理求解方程 $\widetilde F(x,y)=(x,0)$.
\end{Proof}
\endgroup

虽然隐函数定理断言存在唯一解, 但它在唯一性不成立的情形中往往也有用. 考虑方程 $F(x,\lambda)=0$, 其中 $F$ 为从 $X\times\mathbb{R}$ 到 $Y$ 的连续可微函数. 假设对每个 $\lambda$ 都有 $F(0,\lambda)=0$. 若 $D_xF(0,\lambda)$ 既单又满, 则隐函数定理说明 $x=0$ 是 $F(x,\lambda)=0$ 的唯一小解. 值得考察的是使这一假设不成立的那些 $\lambda$ 值. 这引向分岔理论, 关于它已有大量文献. 这里只考虑最简单的分岔情形. 具体地, 假设 $F$ 二次连续可微, $D_xF(0,0)$ 的零空间为由 $x_0$ 张成的一维空间, 其值域的余维为一, 并且 $D_{x\lambda}F(0,0)x_0$ 不在 $D_xF(0,0)$ 的值域中. 现在设 $Z$ 为 $X$ 的一个子空间, 它与 $x_0$ 张成的空间互补; 代入 $x=\epsilon(x_0+z)$, 其中 $z\in Z$. 然后定义 $G(\epsilon,z,\lambda)=\frac{1}{\epsilon}F(\epsilon(x_0+z),\lambda)$, 在极限 $\epsilon=0$ 时显然定义 $G(0,z,\lambda)=D_xF(0,\lambda)(x_0+z)$. 由于假设 $F$ 属于 $C^2$ 类, 可知 $G$ 仍属于 $C^1$ 类. 现在利用隐函数定理求解方程 $G(\epsilon,z,\lambda)$, 将 $z$ 和 $\lambda$ 表为 $\epsilon$ 的函数. 为此, 只需注意导数为 $D_zG(0,0,0)=\left.D_xF(0,0)\right|_Z$ 和 $D_\lambda G(0,0,0)=D_{x\lambda}(0,0)x_0$. 假设 $D_{x\lambda}(0,0)x_0$ 不在 $D_xF(0,0)$ 的值域中, 恰好保证线性化的可逆性. 于是得到方程 $F(x,\lambda)=0$ 的一个非平凡解的分岔分支. 这个非平凡解分支由 ``振幅因子\textquotesingle\textquotesingle\ $\epsilon$ 参数化.
